% pandoc.tex — default.latex de Pandoc con las variables de contenido propias
% de este repositorio. El aspecto lo pone el paquete de `fontfamily`, si lo
% hay; sin él, compila con LaTeX estándar. Divergencias marcadas con DIVERGE.
% Copyright © 2026 Carlos Llamedo (MIT Licence)

$pandoc-templates/document-metadata.latex()$
$pandoc-templates/passoptions.latex()$
% DIVERGE: los nombres de color de xcolor siempre, no solo con colorlinks,
% porque el color de los enlaces lo pone el estilo.
\PassOptionsToPackage{dvipsnames,svgnames,x11names}{xcolor}
\documentclass[
$for(babel-otherlangs)$
  $babel-otherlangs$,
$endfor$
$if(babel-lang)$
  $babel-lang$,
$endif$
$if(fontsize)$
  $fontsize$,
$endif$
$if(papersize)$
  $papersize$paper,
$endif$
$for(classoption)$
  $classoption$$sep$,
$endfor$
]{$documentclass$}
$if(beamerarticle)$
\usepackage{beamerarticle} % needs to be loaded first
$endif$
\usepackage{xcolor}
$if(geometry)$
\usepackage[$for(geometry)$$geometry$$sep$,$endfor$]{geometry}
$endif$
\usepackage{amsmath,amssymb}
$if(cancel)$
\usepackage{cancel}
$endif$
$--
$-- section numbering
$--
$if(numbersections)$
\setcounter{secnumdepth}{$if(secnumdepth)$$secnumdepth$$else$5$endif$}
$else$
\setcounter{secnumdepth}{-\maxdimen} % remove section numbering
$endif$
$pandoc-templates/fonts.latex()$
$pandoc-templates/font-settings.latex()$
$pandoc-templates/common.latex()$
$for(header-includes)$
$header-includes$
$endfor$
$pandoc-templates/after-header-includes.latex()$
% DIVERGE: hypersetup.latex sin el aspecto de los enlaces (colorlinks, colores,
% hidelinks), que es cosa del estilo; solo los metadatos del PDF.
\makeatletter
\@ifundefined{xmpquote}{\newcommand{\xmpquote}[1]{#1}}{}
\makeatother
\hypersetup{
$if(title-meta)$
  pdftitle={$title-meta$},
$endif$
$if(author-meta)$
  pdfauthor={$author-meta$},
$endif$
$if(lang)$
  pdflang={$lang$},
$endif$
$if(subject)$
  pdfsubject={$subject$},
$endif$
$if(keywords)$
  pdfkeywords={$for(keywords)$\xmpquote{$keywords$}$sep$, $endfor$},
$endif$
  pdfcreator={LaTeX via pandoc}}
$if(pdf-trailer-id)$

\ifXeTeX
\special{pdf:trailerid [ $pdf-trailer-id$ ]}
\fi
\ifPDFTeX
\pdftrailerid{}
\pdftrailer{/ID [ $pdf-trailer-id$ ]}
\fi
\ifLuaTeX
\pdfvariable trailerid {[ $pdf-trailer-id$ ]}
\fi
$endif$

$if(title)$
\title{$title$$if(thanks)$\thanks{$thanks$}$endif$}
$endif$
$if(subtitle)$
\usepackage{etoolbox}
\makeatletter
\providecommand{\subtitle}[1]{% add subtitle to \maketitle
  \apptocmd{\@title}{\par {\large #1 \par}}{}{}
}
\makeatother
\subtitle{$subtitle$}
$endif$
% DIVERGE: autores como mapas con name y email; el estilo decide qué hace
% con \authoremail (si no lo define, se ignora). \keywords, si el estilo no
% lo define, es un párrafo con \keywordsname.
\author{$for(author)$$if(author.name)$$author.name$$else$$author$$endif$$sep$ \and $endfor$}
\providecommand{\authoremail}[1]{}
\providecommand{\keywordsname}{Palabras clave}
\providecommand{\keywords}[1]{\par\medskip\noindent\textbf{\keywordsname:} #1}
$for(author)$
$if(author.email)$
\authoremail{$author.email$}
$endif$
$endfor$
\date{$date$}

\begin{document}
$if(has-frontmatter)$
\frontmatter
$endif$
% DIVERGE: portada solo con maketitle; resumen sin depender del título y
% con keywords (\keywords del estilo o el de respaldo).
$if(maketitle)$
\maketitle
$endif$
$if(abstract)$
\begin{abstract}
$abstract$
$if(keywords)$
\keywords{$for(keywords)$$keywords$$sep$, $endfor$}
$endif$
\end{abstract}
$endif$

$for(include-before)$
$include-before$

$endfor$
$if(toc)$
$if(toc-title)$
\renewcommand*\contentsname{$toc-title$}
$endif$
{
$if(colorlinks)$
$if(toccolor)$
\hypersetup{linkcolor=$toccolor$}
$endif$
$endif$
\setcounter{tocdepth}{$toc-depth$}
\tableofcontents
}
$endif$
$if(lof)$
\listoffigures
$endif$
$if(lot)$
\listoftables
$endif$
% DIVERGE: índice de listados.
$if(lol)$
\lstlistoflistings
$endif$
$if(linestretch)$
\setstretch{$linestretch$}
$endif$
$if(has-frontmatter)$
\mainmatter
$endif$
$body$

% DIVERGE: \backmatter siempre que la clase o el estilo lo definan
\ifdefined\backmatter\backmatter\fi
% DIVERGE: variable refs para colocar la bibliografía (véase docs/filtros.md).
$if(refs)$
$refs$
$endif$
$if(nocite-ids)$
\nocite{$for(nocite-ids)$$it$$sep$, $endfor$}
$endif$
$if(natbib)$
$if(bibliography)$
$if(biblio-title)$
$if(has-chapters)$
\renewcommand\bibname{$biblio-title$}
$else$
\renewcommand\refname{$biblio-title$}
$endif$
$endif$
\bibliography{$for(bibliography)$$bibliography$$sep$,$endfor$}

$endif$
$endif$
$if(biblatex)$
\printbibliography$if(biblio-title)$[title=$biblio-title$]$endif$

$endif$
$for(include-after)$
$include-after$

$endfor$
\end{document}
